% ============================================================
% ANNEXE A : GUIDE D'INSTALLATION
% ============================================================

\chapter{Guide d'Installation et de Déploiement}

\section{Prérequis Système}

\subsection{Configuration Minimale}

\begin{table}[H]
\centering
\caption{Configuration minimale requise.}
\renewcommand{\arraystretch}{1.3}
\begin{tabular}{|l|l|}
\hline
\rowcolor{stellantisblue!15}
\textbf{Composant} & \textbf{Exigence minimale} \\
\hline
Système d'exploitation & Ubuntu 20.04+ / Windows 10+ / Raspberry Pi OS \\
\hline
Python & 3.10 ou supérieur \\
\hline
RAM & 4 Go minimum (8 Go recommandé) \\
\hline
Espace disque & 2 Go minimum \\
\hline
Caméra & Compatible UVC (USB Video Class) \\
\hline
GPU (optionnel) & NVIDIA avec CUDA 11.0+ pour l'entraînement \\
\hline
\end{tabular}
\end{table}

\section{Installation sur PC (Développement)}

\begin{lstlisting}[style=pythonstyle, caption={Procédure d'installation complète.}, language={}]
# 1. Cloner le projet
git clone https://github.com/meriemdaifi/dms-project.git
cd dms-project

# 2. Creer un environnement virtuel
python -m venv venv
source venv/bin/activate  # Linux/Mac
# venv\Scripts\activate   # Windows

# 3. Installer les dependances Python
pip install --upgrade pip
pip install -r requirements.txt

# 4. Verifier l'installation
python -c "import cv2; print(f'OpenCV: {cv2.__version__}')"
python -c "import torch; print(f'PyTorch: {torch.__version__}')"

# 5. Lancer les tests unitaires
python -m unittest tests.test_dms -v

# 6. Lancer la simulation principale
python main.py

# 7. Lancer un scenario specifique
python main.py --scenario attentive
python main.py --scenario fatigued
python main.py --scenario distracted
python main.py --scenario mixed
\end{lstlisting}

\section{Installation sur Raspberry Pi (Prototype)}

\begin{lstlisting}[style=pythonstyle, caption={Installation sur Raspberry Pi.}, language={}]
# 1. Mise a jour du systeme
sudo apt update && sudo apt upgrade -y

# 2. Installation des dependances systeme
sudo apt install -y python3-pip python3-venv
sudo apt install -y libopencv-dev python3-opencv
sudo apt install -y libatlas-base-dev libjasper-dev
sudo apt install -y libhdf5-dev libhdf5-serial-dev

# 3. Installation des outils GPIO et CAN
sudo apt install -y python3-rpi.gpio
sudo apt install -y can-utils
sudo pip3 install python-can

# 4. Configuration de l'interface CAN (MCP2515)
# Ajouter dans /boot/config.txt :
# dtoverlay=mcp2515-can0,oscillator=8000000,interrupt=25
# dtoverlay=spi-bcm2835-overlay

# 5. Activation de l'interface CAN
sudo ip link set can0 up type can bitrate 500000

# 6. Cloner et installer le projet
git clone https://github.com/meriemdaifi/dms-project.git
cd dms-project
pip3 install -r requirements_rpi.txt

# 7. Lancer le systeme DMS
python3 main.py --mode robot --camera 0
\end{lstlisting}

\section{Configuration du Module CAN}

\subsection{Câblage MCP2515 -- Raspberry Pi}

\begin{table}[H]
\centering
\caption{Connexions MCP2515 $\leftrightarrow$ Raspberry Pi GPIO.}
\renewcommand{\arraystretch}{1.3}
\begin{tabular}{|l|l|l|}
\hline
\rowcolor{stellantisblue!15}
\textbf{MCP2515} & \textbf{Raspberry Pi} & \textbf{Description} \\
\hline
VCC & Pin 2 (5V) & Alimentation \\
\hline
GND & Pin 6 (GND) & Masse \\
\hline
CS & Pin 24 (GPIO 8 / CE0) & Chip Select SPI \\
\hline
SO & Pin 21 (GPIO 9 / MISO) & SPI Master In Slave Out \\
\hline
SI & Pin 19 (GPIO 10 / MOSI) & SPI Master Out Slave In \\
\hline
SCK & Pin 23 (GPIO 11 / SCLK) & Horloge SPI \\
\hline
INT & Pin 22 (GPIO 25) & Interruption \\
\hline
\end{tabular}
\end{table}

\section{Lancement des Simulations MATLAB}

\begin{lstlisting}[style=matlabstyle, caption={Lancement de la simulation MATLAB.}]
% Ouvrir MATLAB et naviguer vers le dossier du projet
cd('chemin/vers/dms-project/matlab')

% Lancer la simulation complete
dms_simulation

% Ou ouvrir le modele Simulink
open_system('dms_simulink_model')

% Lancer la simulation Simulink
sim('dms_simulink_model', 120)  % 120 secondes
\end{lstlisting}
